\abstract{РЕФЕРАТ}

Объем работы равен \formbytotal{lastpage}{страниц}{е}{ам}{ам}. Работа содержит \formbytotal{figurecnt}{иллюстраци}{ю}{и}{й}, \formbytotal{tablecnt}{таблиц}{у}{ы}{}, \arabic{bibcount} библиографических источников и \formbytotal{числоПлакатов}{лист}{}{а}{ов} графического материала. Количество приложений – 2. Графический материал представлен в приложении А. Фрагменты исходного кода представлены в приложении Б.

Перечень ключевых слов: фитнес, программно-информационная система, клиент, web-приложение, одностраничное приложение, тренировочный план, прогресс, пользователь, интерфейс, профиль, тренировки, питание, социальная лента, DOM, JavaScript, HTML, CSS, REST API, модуль, браузер.

Объектом разработки является клиентская часть программно-информационной системы <<Fitness Life>>, обеспечивающая регистрацию и вход пользователей, работу с тренировочными планами и дневником питания, визуализацию прогресса, взаимодействие с модулями поддержки, аватара и социальной ленты, а также обмен данными с серверной подсистемой через REST API.

Целью выпускной квалификационной работы является создание клиентской подсистемы веб-приложения, обеспечивающей планирование тренировок, отслеживание прогресса, удобный пользовательский интерфейс и синхронизацию персональных данных с серверной частью.

В процессе разработки клиентской части системы была выделена клиент-серверная архитектура, спроектирована модульная структура одностраничного приложения, реализованы экраны аутентификации и профиля, клиентские модули тренировок, питания, психологической поддержки, аватара и социальной ленты, а также механизмы локального хранения состояния и обращения к REST API.

При разработке клиентской части использовались HTML, CSS, JavaScript (модули *.module.js), библиотека Chart.js, Web Storage API (localStorage).

\selectlanguage{english}
\abstract{ABSTRACT}
  
The volume of work is \formbytotal{lastpage}{page}{}{s}{s}. The work contains \formbytotal{figurecnt}{illustration}{}{s}{s}, \formbytotal{tablecnt}{table}{}{s}{s}, \arabic{bibcount} bibliographic sources and \formbytotal{числоПлакатов}{sheet}{}{s}{s} of graphic material. The number of applications is 2. Graphic material is presented in annex A. Source code fragments are presented in annex B.

List of keywords: fitness, information system, client, web application, single-page application, workout plan, progress, user, interface, profile, training, nutrition, social feed, DOM, JavaScript, HTML, CSS, REST API, module, browser.

The object of development is the client part of the Fitness Life information system providing user registration and login, workout planning and nutrition diary, progress visualization, support and avatar modules, social feed interaction, and data exchange with the server subsystem via a REST API.

The purpose of the thesis is to create a client subsystem of a web application that supports workout planning, progress tracking, a convenient user interface, and synchronization of personal data with the server part.

During the development of the client part, a client-server architecture was defined, a modular single-page application structure was designed, authentication and profile screens were implemented, client modules for training, nutrition, psychological support, avatar, and social feed were developed, along with local state storage and REST API access mechanisms.

HTML, CSS, JavaScript (*.module.js modules), the Chart.js library, and the Web Storage API (localStorage) were used in the development of the client part.
\selectlanguage{russian}
